% Chapter 9: Exploratory Morphological Typologies
% Focus: Clustering scientific disciplines into descriptive typologies based on the 11 core embedding-space metrics.
% Document Feeds: docs/TFM_summary.md, docs/reduced_metric_core.md
% Outputs: outputs/04_reduced_metric_core/ (for the 11 metrics), potential clustering tables/plots.
% Citations: Reminders to cite unsupervised clustering methods (e.g. Hierarchical Clustering, K-Means) and typological philosophy.
% IMPORTANT constraints:
%   - Clustering is strictly exploratory and must be based ONLY on the reduced 11 embedding-space metrics, NOT UMAP coordinates.
%   - The clusters must be interpreted as descriptive typologies (conceptual guides), NOT as natural or objective categories of science.
%   - There is NO standalone visualization chapter; UMAP projections and visual density maps are integrated inside the relevant chapters as interpretive support.
% TODO: Run exploratory clustering on the 11-metric profiles and export the resulting cluster members and dendrograms.

\noindent This chapter presents an exploratory clustering analysis of scientific disciplines. Using the eleven core high-dimensional morphological metrics, we group subfields with structurally similar profiles to build descriptive typologies. This unsupervised typology provides a heuristic framework to understand the shared morphological features of modern scientific research.

% TODO: Outline the clustering methodology (e.g., Hierarchical Agglomerative Clustering or K-Means on the 11-metric profiles).
% TODO: Justify clustering directly in the original 768-D derived metric space (excluding UMAP coordinates to prevent distortion).
% TODO: Present the descriptive typologies, analyzing the key morphological traits that define each cluster.
% TODO: Emphasize the epistemological status of clusters as descriptive typologies rather than absolute, objective boundaries.
% TODO: Integrate interpretive UMAP visualizations and density plots directly here to illustrate cluster coherence in 2D.
